% Auto-generated complete chapter from Research/deep-research-report_02.md
\chapter{Auditoría de roles y posicionamiento}
\label{complete:research_055}

\begin{sourcemeta}
\textbf{Fuente:} \href{file:///E:/Laboral/Research/deep-research-report_02.md}{Research/deep-research-report\_02.md}\\
\textbf{Seccion:} research\\
\textbf{Estado:} research\_snapshot\\
\textbf{Rol:} Research / estudio de mercado\\
\textbf{Lineas originales:} 108\\
\textbf{Ultima modificacion:} 2026-06-11 12:16\\
\textbf{Deduplicacion conservadora:} 0 bloques repetidos omitidos; 0 caracteres repetidos no reimpresos.
\end{sourcemeta}

\section*{Resumen de orientacion}
Se evaluaron las 23 familias de roles propuestas en función de la demanda real (2025–2026), encaje con el perfil de Alexander y su objetivo de trabajo remoto desde Colombia. A continuación se resumen las conclusiones clave en tablas y listas.

\section*{Contenido completo depurado}
\addcontentsline{toc}{section}{Contenido completo depurado}




Se evaluaron las 23 familias de roles propuestas en función de la demanda real (2025–2026), encaje con el perfil de Alexander y su objetivo de trabajo remoto desde Colombia. A continuación se resumen las conclusiones clave en tablas y listas.




\subsection{1. Tabla de priorización de roles}




\begin{center}
\scriptsize
\begin{longtable}{>{\raggedright\arraybackslash}p{0.305\linewidth} >{\raggedright\arraybackslash}p{0.305\linewidth} >{\raggedright\arraybackslash}p{0.305\linewidth}}
\toprule
Rol / Posición & Prioridad (Alta/Media/Baja) & Comentarios clave (demanda, requisitos, salario) \\
\midrule
\endfirsthead
\toprule
Rol / Posición & Prioridad (Alta/Media/Baja) & Comentarios clave (demanda, requisitos, salario) \\
\midrule
\endhead
\textbf{Unity Technical Artist} (Artista técnico Unity) & Alta & Rol frecuente en empresas de videojuegos y simulación; portafolio exigente (Shaders, optimización en Unity/URP). Salarios senior \textasciitilde{} \$130–150k. Requiere 3–5+ años de experiencia y dominio de C\#, Python, herramientas 3D. Encaja con TwinSight (Unity+optimización). Remoto posible con EOR/contratación. \\
\textbf{Real-Time 3D Developer / Interactive 3D Developer} & Alta & Roles amplios que incluyen desarrollo interactivo (WebGL, Three.js) y visualización 3D en tiempo real. Demanda creciente en VR/AR y visualización técnica. TwinSight es ejemplo directo. Salarios muy variables; startups WebGL pueden pagar hasta \textasciitilde{}\$125k. Requiere Unity/WebGL y optimización. Buen encaje con perfil (C\#, WebGL, UX Toolkit). \\
\textbf{Technical Artist (general)} (Artista técnico) & Media-Alta & Demanda sólida en estudios de juegos y simulación (pipeline entre arte y dev). Similares requisitos a \emph{Unity TA} pero más genérico (puede incluir Unreal, Houdini, etc.). Salarios promedio \textasciitilde{} \$138k EE. UU. (higher than entry targets). Requiere experiencia en herramientas 3D (Photoshop, Maya, Houdini) y scripting. Buen encaje por know-how gráfico, pero muy competitivo. \\
\textbf{XR/AR/VR Developer} (Desarrollador XR/AR/VR) & Media & Roles emergentes en empresas de realidad extendida (Meta, Google, startups). Requiere experiencia en Unity XR, C\#, diseño de experiencias 3D. Ofertas Sr. XR dev \textasciitilde{} \$70–100k; puestos remotos (\textasciitilde{}\$90k) existen. Perfil puede pivotar hacia AR/VR usando TwinSight si amplía skills. Idiomas hablados pueden requerirse (inglés). Demanda creciente en simulación/educación/industrial. \\
\textbf{Technical Visualization / Digital Twin Developer} & Media & Orientado a industrias (manufactura, energía, defensa) que necesitan crear gemelos digitales o visualizaciones técnicas. Roles pocos en cantidad, pero suelen requerir Unity o motores 3D (como en Cincinnati ejemplo). Salarios en EE. UU. \textasciitilde{}\$100–160k. Gran encaje para TwinSight (drone CAD→Unity). Remoto posible (muchas consultoras), pero competido. \\
\textbf{Unity / WebGL Developer} (general) & Media & Amplia demanda global en videojuegos y aplicaciones 3D web. Pueden incluir desde front-end 3D hasta AR/VR. Ejemplo: Unity developer industrial (HMI) con Unity. Salario promedio WebGL \textasciitilde{} \$125k en startups, pero en promedio USA \textasciitilde{}\$21/h. Remoto LATAM viable; ejemplos pagaron \$5–6k mensuales (Somewhere). Perfil coincide en parte con TwinSight/WebGL; requiere marketing de portafolio web. \\
\textbf{Tools/Pipeline Developer} (Desarrollador de herramientas) & Media & Enfocado en scripting y automatización (blender, pipeline, ARA). Los roles puros son raros; a veces aparece en estudios grandes. Salario medio un Python dev (\textasciitilde{}\$80k–120k). TwinSight/ARA evidencia interés, pero necesitan traducir a aplicación de producto (p.ej. editor de Unity o herramienta CAD). Competencia moderada, requiere demostrar proyectos útiles (ARA framework). Remoto factible. \\
\textbf{3D Product Visualization Developer} (Gráficos 3D industria) & Media & Empresas de e-commerce o manufactura crean vistas 3D de productos. Requiere Blender/Unity para presentaciones interactivas. Demanda menor pero creciente; salarios variables (p.ej. industria automotriz/healthcare 80–120k). Encaje alto con Blender Portrait, pero hay mucha competencia de artistas 3D y demanda de UI/UX. \\
\textbf{Simulation / Engineering Developer} & Media & Incluye roles de simulación industrial o entrenamiento (e.g. EEUU, Defensa). A menudo Unity/Unreal con C\#. Ofertas muestran desarrolladores 3D XR para el ejército (\textasciitilde{}\$70–90k). Requiere background técnico (física/matemática) que Alexander tiene. Permite aprovechar skills físicos. Remoto limitado (muchos on-site en Defensa), aunque consultoras contratan. \\
\textbf{Shader Technical Artist} & Media & Rol especializado en shaders, efectos gráficos (Unity/Unreal). Demanda específica en producción AAA o VFX. Salarios senior altos (\$100k+). Requiere portafolio de shaders (Shader Graph, HLSL) y conocimiento de URP. Perfil Alexander tiene bases (Shader Graph, HLSL básico). Puede usarse en portfolios. \\
\textbf{Full-Stack Developer} & Baja & Demanda general en software (ya sea web, apps). Pero se aleja del eje 3D/Unity. Aunque hay muchas vacantes globales y remotas (React/Python/DB), sería desviar la carrera. Salarios \textasciitilde{} \$60–120k. Se considera ruta de contingencia si no se consigue otro rol, pero no alineado con portafolio técnico-artístico. \\
\textbf{3D Artist / Generalist (Blender)} & Baja & Rol creativo puro, muy competitivo. Perfil de Alexander es técnico (no enfocarse en arte estático). Además, portafolios de arte 3D saturados (mostrados en ArtStation). Salario medio menor (\$40–60k). Se puede incluir un Blender Portrait de respaldo, pero no debe ser el foco principal. \\
\textbf{AI Tools / Python Developer} & Baja & Roles de inteligencia artificial e integración de LLM (ej. \emph{AI Tools Developer}, \emph{LLM Developer}). Creciente demanda (data science, ML), pero requiere experiencia sólida en ML. ARA/LLM aquí es un prototipo interno; no es un producto terminado. Ofertas de Python genérico son comunes (50–120k), pero el foco de Alexander es 3D más que IA. Puede usarse de forma secundaria (herramientas de pipeline), pero no es primaria en corto plazo. \\
\textbf{CAD-to-Realtime Tech Artist} & Baja & Extremadamente nicho. No aparecen ofertas específicas de “CAD-to-Realtime” como título. Implica industria CAD (autocad, NX) + optimización 3D. Podría encajar en \emph{Technical Visualization} o \emph{Digital Twin}, pero no parece un rol autónomo. Más bien, es una habilidad dentro de dev de visualización técnica. \\
\bottomrule
\end{longtable}
\end{center}




\begin{quote}
\textbf{Fuentes:} Se consultaron portales de empleo y comparadores salariales. Por ejemplo, ZipRecruiter indica un salario medio \textasciitilde{}\textbackslash{}\$138k/año para \emph{Technical Artist} en EE.UU.. LinkedIn/Indeed muestran \emph{Unity Developer} LATAM ofreciendo \textbackslash{}\$5–6k/mes. El sector Web3 startup reporta \textbackslash{}\$125k/yr para \emph{WebGL Developer}. Para XR dev remotos se ven rangos \textbackslash{}\$70–100k. Los datos reflejan el mercado a mediados de 2026 y los ajustamos a la realidad de Colombia (contratos B2B o EOR).
\end{quote}




\subsection{2. Roles secundarios y rutas de contingencia}




\begin{itemize}
\item \textbf{Technical Artist general (no Unity):} Similar al rol principal pero ampliando a otras tecnologías (Unreal, Maya, Houdini). Como ruta secundaria, útil si no hay vacantes Unity específicas. Exige experiencia similar.
\item \textbf{XR/Simulation Developer:} Desarrollo de aplicaciones VR/AR o simulaciones técnicas. Puede apalancar el mismo skillset (3D, C\#) y la formación en física, pero suele requerir manejo de dispositivos específicos (oculus, hololens) y aún es competitivo.
\item \textbf{Interactive/WebGL Developer:} Especializado en experiencias 3D Web (Three.js, WebAssembly). Aprovecha la tesis TwinSight (WebGL) y la habilidad Python/C\#. No tan bien pagado como otros, pero buena puerta de entrada.
\item \textbf{Tools/Pipeline Developer:} Para roles internos de automatización (Unity plugins, herramientas de animación). Utiliza ARA/AI skills. Mercado limitado, a veces parte del equipo de tech art. Salario medio.
\item \textbf{Game Technical Artist:} Centrado en videojuegos AAA. Muy competitivo para un perfil \emph{junior}. Buena fuente de motivación y apta si se mejora portfolio con demos de juegos. Pero en 3–6 meses quizá poco realista llegar a ofrecer.
\end{itemize}




\textbf{Rutas de contingencia (fallback):} full-stack (desarrollador web/C\#) como último recurso; proyectos freelance (Upwork, Freelancer) en Unity/WebGL; apoyo en diseño 3D de producto (bajos pagos pero refuerza portfolio); continuar freelance local mientras se busca remoto.




\subsection{3. Posicionamiento principal recomendado}




Basado en la tabla anterior, se sugiere un \textbf{posicionamiento único centrado en 3D interactivo y visualización técnica real-time}. Por ejemplo:




\begin{quote}
“\emph{Desarrollador 3D en tiempo real y Artista Técnico especializado en Unity, C\# y WebGL, enfocado en visualización técnica interactiva (simulaciones, gemelos digitales, visualización CAD‑a‑tiempo real).}”
\end{quote}




Este enunciado destaca Unity/C\#, WebGL y la idea de visualización técnica (TwinSight) sin limitarse a videojuegos. Enfatiza habilidades fuertes (optimización, herramientas 3D) y el beneficio clave (“visualización técnica interactiva”).




\subsection{4. Rutas secundarias de posicionamiento}




\begin{itemize}
\item \textbf{Interactive 3D/WebGL Developer:} Enfocarse en portales, sitios o experiencias web 3D. LinkedIn y GitHub deben reflejar proyectos con Three.js/WebGL.
\item \textbf{XR/AR/VR Developer:} Resaltar interés y base física/simulación; incluir cualquier demostración experimental en AR/VR.
\item \textbf{Digital Twin/Simulación:} Mencionar perspectiva industrial; destacar reducción de polígonos CAD→3D (TwinSight).
\item \textbf{Python/IA Tooling (ARA):} Como ruta auxiliar, puede presentar el framework ARA como muestra de habilidad en Python y LLM, pero aclarando que es prototipo.
\end{itemize}




\subsection{5. Rutas de contingencia}




En caso de dificultades iniciales, considerar:


\begin{itemize}
\item Proyectos freelance cortos en Unity/WebGL (contratación por horas o app móvil simple).
\item Roles de desarrollador general (C\#, .NET) en empresas de software (salario más bajo, usa portfolio como backend).
\item Mejorar portafolio con cursos prácticos (Blender, shader).
\item Preparación de certificaciones (Unity, etc.) para ganar credibilidad sin gastar años.
\end{itemize}




\subsection{6. Keywords ATS sugeridas por rol}




\_Para optimizar CV/LinkedIn/GitHub ante sistemas de búsqueda:\_




\begin{itemize}
\item \textbf{Unity Technical Artist:} \texttt{Unity}, \texttt{C\#}, \texttt{Shader Graph}, \texttt{URP}, \texttt{optimization}, \texttt{pipeline}, \texttt{content creation tools (Photoshop/Maya/Houdini)}, \texttt{technical art}, \texttt{real-time rendering}, \texttt{VR}, \texttt{WebGL}.
\item \textbf{Technical Artist (general):} \texttt{Technical Artist}, \texttt{game pipeline}, \texttt{VFX}, \texttt{Houdini}, \texttt{Substance}, \texttt{tool development}, \texttt{Python}, \texttt{C++}, \texttt{3D optimization}.
\item \textbf{Real-Time 3D/Interactive Dev:} \texttt{Unity}, \texttt{WebGL}, \texttt{Three.js}, \texttt{JavaScript}, \texttt{graphics performance}, \texttt{C\#}, \texttt{interactive UI}, \texttt{UI Toolkit}, \texttt{runtime systems}, \texttt{mobile optimization}.
\item \textbf{Unity/WebGL Developer:} \texttt{Unity3D}, \texttt{WebGL}, \texttt{HTML5}, \texttt{JavaScript}, \texttt{CI/CD}, \texttt{Addressables}, \texttt{profiling}, \texttt{Async/Await}, \texttt{performance}.
\item \textbf{Technical Visualization/Digital Twin:} \texttt{Digital Twin}, \texttt{CAD}, \texttt{PiXYZ}, \texttt{Blender}, \texttt{tessellation}, \texttt{simulation}, \texttt{BIM}, \texttt{MAYK} (Manufacturing, Aerospace, Militar, etc.), \texttt{STL}.
\item \textbf{XR/AR/VR Developer:} \texttt{XR}, \texttt{AR}, \texttt{VR}, \texttt{HoloLens}, \texttt{ARCore/ARKit}, \texttt{3D interaction}, \texttt{spatial mapping}, \texttt{Unity XR}, \texttt{Oculus}, \texttt{Meta}.
\item \textbf{Tools/Pipeline Dev:} \texttt{Python}, \texttt{LLM}, \texttt{LangChain}, \texttt{FastAPI}, \texttt{Blender Add-on}, \texttt{CI/CD}, \texttt{automation}, \texttt{editor scripting}, \texttt{software tooling}.
\item \textbf{Shader Technical Artist:} \texttt{HLSL}, \texttt{ShaderLab}, \texttt{URP}, \texttt{High Definition Render Pipeline (HDRP)}, \texttt{Shader Graph}, \texttt{GPU Compute}, \texttt{post-processing}.
\item \textbf{Python/AI Tools Developer:} \texttt{Python}, \texttt{LLM}, \texttt{AI}, \texttt{Machine Learning}, \texttt{LangGraph}, \texttt{agent}, \texttt{automation}, \texttt{TensorFlow}, \texttt{PyTorch}, \texttt{API}.
\item \textbf{Full-Stack (fallback):} \texttt{React}, \texttt{Node.js}, \texttt{FastAPI}, \texttt{PostgreSQL}, \texttt{REST}, \texttt{Docker}, \texttt{AWS}, \texttt{Git}.
\item \textbf{3D Artist / Generalist:} \texttt{Blender}, \texttt{3D modeling}, \texttt{texturing}, \texttt{rigging}, \texttt{UV unwrapping}, \texttt{PBR materials}, \texttt{lighting}.
\end{itemize}




\subsection{7. Gaps de evidencia por rol}




\begin{itemize}
\item \textbf{Roles CAD‑to‑Realtime y CAD‑Centric:} No hay muchos anuncios directos; suelen encajar en \emph{Technical Visualization}. Investigar empresas de ingeniería/defensa (p.ej. Unity for Industry, Dassault, Siemens NX).
\item \textbf{XR/AR:} Falta dato concreto de demanda para entrada JR; solo posiciones senior visibles (Google, Meta). Investigación de nicho VR é importante.
\item \textbf{Herramientas/AI (ARA, AI-Tools):} Roles específicos casi no aparecen en portales. ARA es un prototipo interno; debería definirse si se hará open-source. Riesgo de canibalizar el branding si se muestra como “proyecto principal” sin mercado directo.
\item \textbf{Full-Stack:} Evidencia abundante, pero no alineada con meta principal; uso como evidencia de flexibilidad si se convierte en opción secundaria.
\item \textbf{Cursos/Rebelway:} No hemos evaluado aún si presentar cursos avanzados (Blender, CG) en CV/LinkedIn. Debe validarse cuáles aportan por sí mismos sin certificación.
\item \textbf{Idioma Alemán:} No se encontró evidencia de rol inmediato que exija alemán; generalmente Inglés es suficiente en remote global, alemán solo en roles locales presenciales.
\end{itemize}




\subsection{8. Contradicciones estratégicas a evitar}




\begin{itemize}
\item \textbf{Documentación de proyectos:} No confundir el contenido de tesis. La tesis final es \emph{TwinSight X500} (dron) – \textbf{no} el ARA Framework. En CV/LinkedIn, se debe corregir para evitar malentendidos.
\item \textbf{Citizenship/Movilidad:} No afirmar derechos de trabajo en Europa sin haber obtenido ciudadanía o visa. Enfatizar “remoto desde Colombia” y mencionar la ciudadanía portuguesa como plan futuro, no actual.
\item \textbf{Senioridad exagerada:} Evitar términos como “Senior” o “Expert” en CV si no se cuenta experiencia formal (ACTA mostró rol técnica sin empleo fijo). Mejor usar “Junior-Mid” y enfocar en proyectos concretos.
\item \textbf{Enfoque de rol mixto:} No mezclar en un mismo perfil demasiadas áreas sin conexión (p.ej. “Full-Stack” + “Artista técnico” + “Data”). Esto diluye la estrategia. Enfocar en “3D en tiempo real” prioritariamente.
\item \textbf{Exagerar aporte de IA:} En entrevistas/portafolio, aclarar que ARA fue un proyecto personal prototipo, donde se usó IA como herramienta, pero \textbf{no} autopresentarlo como producto comercial. La responsabilidad técnica final fue del candidato.
\item \textbf{Expectativas salariales altas prematuramente:} Evitar anunciar objetivos de \$6k/mes en 6 meses sin experiencia comprobable. Mejor posicionarse para ofertas de \$3k–\$4k/mes inicialmente (viajar hasta \$6k con mejoras en skills).
\item \textbf{Uso de certificaciones no oficiales:} No dar demasiada importancia a cursos sin certificación (CG Cookie, Alive!), mencionarlos como “formación práctica” si es necesario, pero no como credenciales.
\end{itemize}




\subsection{9. Recomendación final: principales familias de rol (3–5) a atacar}




\begin{enumerate}
\item \textbf{Artista Técnico/Desarrollador 3D en tiempo real (Unity)}: Este es el rol principal. Incluye posiciones como \emph{Unity Technical Artist}, \emph{Technical Artist}, \emph{Real-Time 3D Developer}. Enfatizar su experiencia en proyectos Unity (TwinSight), optimización de activos CAD→Unity, y habilidades en C\#/URP/Shader Graph. Apuntar a empresas de simulación, digital twins e incluso videojuegos. Salario objetivo: \$2–3k mensuales para empezar, escalando rápido con experiencia (target a \$6k en 1–2 años acorde al mercado).
\end{enumerate}




\begin{enumerate}
\item \textbf{Desarrollador XR/Simulación/Digital Twin}: Roles en realidad extendida y simulación industrial (posiciones de \emph{XR Developer}, \emph{Simulation Developer}, \emph{Digital Twin Developer}). Utiliza la formación en electrónica/física y Unity para crear simulaciones inmersivas. Enfocar en empresas que trabajen con entrenamiento militar, salud, manufactura o ciudades inteligentes. Ejemplos: desarrollos AR en Unity, gemelos digitales en Unity/Houdini, apps VR técnicas. Requieren aprendizaje de SDKs XR (ARCore, etc.). Su potencial salarial puede llegar a \$70–90k (EE. UU.). 
\end{enumerate}




\begin{enumerate}
\item \textbf{Desarrollador WebGL/Interactivo 3D}: Familias como \emph{Interactive 3D Developer}, \emph{Unity WebGL Developer} y \emph{Frontend 3D Visualization Developer}. Aquí aprovecha la tesis TwinSight (WebGL) y habilidades de frontend. Empresas de tecnologías Web3, marketing inmersivo y educación ofrecen roles remotos. Aunque el salario inicial es más bajo (\$40–60k), se sube rápido con portafolio y experiencia. Destacar proyectos personales en WebGL (TwinSight) en GitHub/LinkedIn.
\end{enumerate}




\begin{enumerate}
\item \textbf{Herramientas de Pipeline / Automatización (AI)}: Como ruta secundaria ambiciosa, presentar ARA Framework como ejemplo de destrezas en Python y LLM. Roles en empresas que integran LLM en flujos de trabajo (p.ej. GenAI, data-mining) pueden interesar. Sin embargo, \textbf{no priorizar} esta ruta ahora; usarla más como diferenciador suplementario que como foco. 
\end{enumerate}




\begin{enumerate}
\item \textbf{Fallback – Full-Stack / Generalista}: Solo en caso de necesidad inmediata. Posee conocimientos de Python/FastAPI/React pero el mercado es altamente competido. Se sugiere como opción secundaria baja prioridad, manteniendo siempre el foco en roles 3D/Unity arriba. 
\end{enumerate}




\textbf{Resumen:} Priorizar propuestas alineadas a \emph{Unity/C\# en 3D interactivo y visualización técnica}, evitando distraerse en roles horizontales sin conexión al eje. Seguir estas familias de rol, con la estrategia de portafolio y networking correspondiente, maximizará las probabilidades de conseguir empleo remoto bien remunerado en 6–12 meses.




\textbf{Fuentes:} Datos obtenidos de ofertas reales y reportes de mercado (ZipRecruiter, Indeed, Wellfound). Ejemplos: Oferta \emph{Technical Artist} en Activision (Carlsbad) pedía 5+ años y pagaba \textasciitilde{}\$130–156k; posiciones Unity en LATAM ofrecen \textbackslash{}\$5–6k mensuales; mercado WebGL Startups muestra salarios promedio \textbackslash{}\$125k. Estos respaldan las estimaciones anteriores.



